\label{sec:literature}
\tikzset{
  litbox/.style={draw,rounded corners=2pt,fill=gray!12,align=center,font=\footnotesize,minimum height=7mm,inner sep=3pt},
  litio/.style={draw,rounded corners=6pt,fill=white,align=center,font=\footnotesize,minimum height=7mm,inner sep=3pt},
  litarr/.style={-{Stealth[length=2.0mm]},thick},
  litlab/.style={font=\scriptsize,inner sep=1.5pt,align=center}
}
\subsection{Handwriting recognition and writer identification}
\label{sec:lit-background}
\label{sec:lit-hwr}
\subsubsection{Offline recognition and available datasets}
The ability of a computer or system to identify written text after being trained on large amounts of data is known as HWR (Handwriting Recognition) or HTR (Handwritten Text Recognition)~\cite[p.~39]{chen2023}. One type of handwritten text recognition is offline recognition, which means the data being identified by the system is static and captured post-completion of the page of handwritten text. This is ideal for identifying the text from a fixed image of handwritten content by utilising image processing techniques and algorithms to identify the text that was written~\cite[p.~268]{ma2023}. The other type is online recognition, which is a more dynamic process that uses information recorded during the creation of the handwriting. The addition of a time variable includes data about the order of strokes as well as pen-up or pen-down movements~\cite[p.~18]{graves2013}.

This project uses offline recognition because its inputs are photographed or scanned pages. The text and writer must therefore be identified from visible ink geometry. The absence of recorded stroke order also influences how the handwriting can later be reproduced: a path recovered from an image is a geometric estimate, rather than a recording of the writer's movements.

Multiple datasets have been used for training handwritten text recognition modules. The EMNIST (Extended MNIST) By\_Class dataset consists of 814\,255 handwritten characters and numbers, each being 28 by 28 pixels in size~\cite[p.~40]{chen2023}. There is only a single character per image, which suits isolated-character classification but does not itself supply examples of connected handwritten words. The IAM handwriting dataset provides complete handwritten English pages, lines and words with their transcriptions and writer metadata~\cite{marti2002}. Hemanth et al. report the expanded database as 5\,685 labelled sentences and 115\,320 labelled words from 657 writers~\cite[p.~2460]{hemanth2022}.

The dataset influences which variations the classifier encounters and which labels are available. For this project, line images and their transcripts support recognition without manually labelling every character boundary. Writer-labelled pages also support author classification and extraction of individual style profiles. IAM supplies public handwriting variation, while the personal dataset represents the project's photographed input conditions. A small isolated-character dataset alone would not address both requirements.

\subsubsection{Image conditioning and text separation}
The image conditioning of handwriting pages includes techniques performed on the handwritten text before it is passed into the classification and feature-recognition algorithms. This accounts for dimensional adjustments, removal of inconsistencies and normalisation of the data, to make the images more appropriate for text separation and classification~\cite[pp.~2458--2460]{hemanth2022}. Normalisation reduces the influence of unwanted variation; examples include min--max scaling and z-score normalisation~\cite[pp.~40--41]{chen2023}.

Binarisation differentiates the background page from the ink of the written text, computed by pixel. A global threshold taken from the grey-level histogram uses one decision level for the image~\cite{otsu1979}. A local threshold instead uses statistics from a window around each pixel and can accommodate uneven illumination~\cite[pp.~2--3]{bradley2007}. Deskewing corrects page rotation, while deslanting corrects the lean of individual characters by applying a horizontal shear. These are different transformations: aligning text lines does not require removing the writer's characteristic slant~\cite{postl1986,vinciarelli2001}.

For the camera input in this project, illumination correction and deskewing are needed before the ink is separated into lines. Slant must remain available as a writer-specific measurement; it can be removed for glyph storage and reapplied when the text is synthesised. The purpose is to suppress capture variation while retaining the information needed for reproduction.

Text separation is the process of breaking up a page of handwritten text into smaller segments or blocks that can be used in feature extraction and neural-network classifiers. Existing methods separate text into lines, words or individual characters. Hemanth et al. obtain lines and words through filtering, dilation and bounding boxes~\cite[pp.~2459--2460]{hemanth2022}. When lines are mostly straight and separated by blank rows, boundaries can be found from minima in the row ink profile. Where ascenders and descenders overlap, a boundary that bends between lines is more suitable. Minimum-cost seam methods provide a way to trace such curved boundaries~\cite[p.~3]{avidan2007}.

The project therefore separates whole lines before recognition. This avoids requiring a reliable cut between every pair of connected letters. Character boundaries are still needed to construct a glyph library, but there the known transcript and a trained line reader can constrain the extraction. Line segmentation and glyph extraction consequently solve different problems, even though both divide an image.

\subsubsection{Neural networks and feature extraction}
Many different types of neural-network models have been implemented for handwriting recognition, including artificial neural networks (ANNs), convolutional neural networks (CNNs), recurrent neural networks (RNNs) and combinations of them. ANNs consist of interconnected neurons, weights, hidden layers and biases that learn a mapping from the input to an output. CNNs use local connections and shared kernels to extract spatial patterns with fewer independent parameters than an otherwise comparable fully connected image layer~\cite[pp.~6--8]{lecun1998}. RNNs maintain a state across a sequence, allowing a prediction to depend on neighbouring evidence~\cite[p.~269]{ma2023}.

The convolutional layers of a CNN perform feature extraction by sliding a convolutional kernel over the image. Each local input value is multiplied by a corresponding kernel value, the products are summed, and a bias and activation function produce an output feature value. Reusing the kernel allows the same pattern to be detected at different locations~\cite[p.~6]{lecun1998}. For unit stride and no padding, an $N$ by $N$ input with an $f$ by $f$ kernel returns an $(N-f+1)$ by $(N-f+1)$ feature map. Pooling then reduces spatial resolution and retains selected responses~\cite[pp.~6--7]{lecun1998}. Figure~\ref{fig:lit-conv} illustrates these operations.

\begin{figure}[!htbp]
\centering
\begin{tikzpicture}[x=1cm,y=1cm,font=\footnotesize]

\draw[step=0.26,gray!50,thin] (0,0) grid (1.56,1.56);
\draw[thick] (0,0) rectangle (1.56,1.56);
\fill[gray!40,opacity=0.75] (0.26,0.78) rectangle (1.04,1.56);
\draw[very thick] (0.26,0.78) rectangle (1.04,1.56);
\node[litlab,below] at (0.78,-0.10) {input, $N\times N$};
\node[litlab,above] at (0.78,1.62) {$f\times f$ kernel slides};

\draw[litarr] (1.75,0.78) -- (2.60,0.78);
\node[litlab,above] at (2.17,0.86) {multiply,};
\node[litlab,below] at (2.17,0.72) {add bias};

\draw[step=0.26,gray!50,thin] (2.80,0.26) grid (3.84,1.30);
\draw[thick] (2.80,0.26) rectangle (3.84,1.30);
\fill[gray!40,opacity=0.75] (2.80,1.04) rectangle (3.06,1.30);
\node[litlab,below] at (3.32,0.03) {feature map,};
\node[litlab,below] at (3.32,-0.22) {$(N-f+1)\times(N-f+1)$};
\node[litlab,above] at (3.32,1.49) {activation applied};

\draw[litarr] (4.03,0.78) -- (4.88,0.78);
\node[litlab,above] at (4.45,0.86) {max-pool};
\node[litlab,below] at (4.45,0.72) {$2\times2$};

\draw[step=0.26,gray!50,thin] (5.08,0.52) grid (5.60,1.04);
\draw[thick] (5.08,0.52) rectangle (5.60,1.04);
\node[litlab,below] at (5.34,0.29) {downsampled};
\node[litlab,below] at (5.34,0.04) {feature map};

\draw[litarr,dashed] (5.79,0.78) -- (6.55,0.78);
\node[litlab,right,align=left] at (6.60,0.78) {repeated, each stage\\ holding more feature\\ maps than the last};
\end{tikzpicture}
\caption{One convolutional layer and one max-pooling layer. The kernel slides over every position of the input, so an $N\times N$ input gives an $(N-f+1)\times(N-f+1)$ feature map, which pooling then downsamples~\cite[pp.~6--8]{lecun1998}. This size expression assumes unit stride and no padding.}
\label{fig:lit-conv}
\end{figure}


The LeNet-5 model is an established example of this structure. It uses a $32\times32$ input and seven trainable layers, including convolution, subsampling and fully connected stages; its original output uses radial-basis units~\cite[pp.~7--8]{lecun1998}. A deeper alternative is a residual network, whose skip connections provide additional routes for information and gradients through the network~\cite[p.~270]{ma2023}. These examples show the choice between compact feature extraction and a more complex model; increasing depth does not guarantee an improvement for every dataset or processor.

A combination of CNNs and recurrent layers allows an entire line to be read without word or character separation. The convolutional layers convert the cropped line into feature columns, and a bidirectional recurrent layer reads the sequence in both directions, providing context from before and after each position~\cite[pp.~4--5]{kizilirmak2023}. In this offline setting, the sequence axis is the horizontal position in the image, rather than the measured time at which a pen moved. The full line crop is available before classification. Related convolutional--recurrent recognition systems use the same distinction between image features and sequence prediction~\cite[pp.~2--3]{shi2017}.

Long sequences make simple recurrent networks difficult to train because gradients can diminish or grow through repeated updates. Long short-term memory (LSTM) uses an additive memory state and gates to control what is retained, replaced and exposed~\cite[pp.~4--5]{graves2013}. The convolutional feature columns usually outnumber the transcript characters, so their alignment is initially unknown. Connectionist temporal classification (CTC) adds a blank symbol and sums the probabilities of the paths consistent with the transcript. Repeated adjacent symbols are collapsed before blanks are removed~\cite[pp.~2--4]{graves2006}. Figure~\ref{fig:lit-line} shows the resulting recognition path.

\begin{figure}[!htbp]
\centering
\begin{tikzpicture}[x=1cm,y=1cm,font=\footnotesize,node distance=4.5mm]
\node[litio,text width=1.5cm] (img) {line image};
\node[litbox,right=of img,text width=1.7cm] (cnn) {convolutional\\layers};
\node[litbox,right=of cnn,text width=1.6cm] (cols) {$T$ feature\\columns};
\node[litbox,right=of cols,text width=1.9cm] (rnn) {recurrent layer,\\both directions};
\node[litbox,right=of rnn,text width=1.6cm] (pc) {symbol per\\column};
\node[litio,right=of pc,text width=1.1cm] (out) {\texttt{the}};
\foreach \a/\b in {img/cnn,cnn/cols,cols/rnn,rnn/pc,pc/out}{\draw[litarr] (\a) -- (\b);}

\draw[litarr] ([yshift=3.2mm]rnn.north west) -- ([yshift=3.2mm]rnn.north east);
\draw[litarr] ([yshift=6.0mm]rnn.north east) -- ([yshift=6.0mm]rnn.north west);

\node[litlab,below=2.0mm of pc,font=\ttfamily\scriptsize] (sym) {-tt-h-ee};
\node[litlab,below=1.6mm of sym] {collapse repeats,\\drop the blanks};
\node[litlab,below=2.0mm of cols,text width=2.5cm] {$T$ far exceeds the\\number of characters,\\so the alignment is\\not known};
\end{tikzpicture}
\caption{Reading a whole text line without cutting out characters, after the arrangement of Kizilirmak and Yanikoglu~\cite[pp.~4--5]{kizilirmak2023}. The convolutional layers produce $T$ feature columns, a recurrent layer reads them in both directions, and a symbol including a blank is emitted per column; collapsing repeats and removing the blanks gives the transcription.}
\label{fig:lit-line}
\end{figure}


This combination is used for the project's reader because the training labels describe complete lines. The CNN extracts stroke patterns, the bidirectional LSTM retains their horizontal context, and CTC trains the network without supplied character positions. For author identification, the same convolutional architecture is useful, but its features must describe the writer across a line rather than decode an ordered transcript. The two tasks therefore require different output heads and separately trained weights.

\subsubsection{Training and assessment}
The way networks are trained corresponds to the architecture they are built from. Batch normalisation subtracts the batch mean and divides by the standard deviation, with a stabilising constant and learned scale and offset~\cite[p.~4]{ioffe2015}. Dropout reduces overfitting by removing a random subset of a layer's outputs during training, limiting dependence on individual features~\cite[pp.~1931, 1934]{srivastava2014}. Data augmentation introduces small rotations, shears or smooth distortions while retaining the transcription~\cite{wigington2017}. These methods address different sources of variation and do not replace separate validation data.

Recognition is commonly assessed through character error rate (CER) or word error rate (WER), using insertions, deletions and substitutions relative to a reference transcript~\cite[p.~2]{graves2006}. An error in one character can spoil a whole word, but word edit error is not simply a count of words at matching positions. The metric and its reference population must remain explicit when comparing systems.

For this project, numerical forward passes, backward passes and optimiser updates make the learning calculations inspectable. A corresponding library implementation supports training and comparison, with stored arrays used for NumPy inference on the ODROID. The report distinguishes the selected library-trained weights from the independently trained numerical models. Character edits, writer decisions and physical drawing measurements are evaluated separately because they measure different functions.

\subsubsection{Techniques for writer identification}
There are multiple approaches to identifying the author of a written piece of text. One uses measured, text-independent features such as slant, height, stroke width, spacing and connectivity. Another describes the distribution of edge directions and recurring letter fragments, giving a texture and shape description of the writing~\cite[p.~701]{bulacu2007}. Alternatively, a neural network learns discriminative features from examples labelled by writer~\cite{he2020}. Such a classifier identifies a writer among its trained classes; it does not automatically identify an unfamiliar person.

This project needs an author classifier to select and assess the stored writer styles, while explicit geometry measurements are needed to reconstruct their appearance. Ink-trace width can itself contain writer information~\cite[p.~162]{brink2012}, but the gantry uses its own pen. Stroke normalisation is therefore applied to reduce reliance on that cue and direct the classifier towards shape. Writer identity and stylistic fidelity remain separate measures: recognising an author does not establish that every letter or spacing feature has been reproduced accurately.

\subsection{Handwriting synthesis in a writer's style}
\label{sec:lit-synth}
Before handwritten text can be reproduced through hardware, the desired text must be created in the selected author's style. Four approaches illustrate the trade-offs between source data, learned variation and the form of the output.

\textbf{Recurrent stroke-level synthesis.} A recurrent neural network predicts the next pen movement through small coordinate offsets and an end-of-stroke indicator. It models a distribution of possible movements because no letter is written identically every time~\cite[pp.~18--20]{graves2013}. Its output is naturally suited to a plotter, but training uses online recordings. This project's offline pages contain no measured stroke order, so this approach would require a different source dataset and an additional generative model.

\textbf{Adversarial image synthesis.} A generative adversarial network trains a generator and discriminator together: one generates handwriting, while the other helps distinguish generated samples from real examples. GANwriting also constrains content and writer style to produce text images~\cite[pp.~3--7]{kang2020}. Once trained, the generator produces an image in a forward pass. This is a relevant alternative for learning contextual letterforms, but a raster output still needs conversion into the pen paths required by this project.

\textbf{Diffusion-based synthesis.} This method learns to reverse a noise-corruption process and generates handwriting through successive denoising steps. One-DM uses one reference style image at generation time, while its training uses a much larger handwriting corpus~\cite[pp.~3--5, 9]{dai2024}. A small reference requirement must therefore be distinguished from the data and computation required to train the model. For this project, diffusion would add both a separate training task and image-to-path recovery, rather than directly supplying the motion representation.

\textbf{Glyph-library synthesis.} This method assembles words from shapes or glyphs extracted from existing handwritten pages. Haines et al. store labelled forms and select and join them using information about their shape, context and spacing~\cite[pp.~4--8]{haines2016}. Their standard data collection uses approximately four A4 pages per writer~\cite[p.~4]{haines2016}. This suits the project's small personal collections and allows explicit control of size, slant and spacing. The main limitation is that a finite collection of character shapes and joining rules only approximates the way connected handwriting changes with its neighbours.

Generated handwriting can be assessed through geometric comparisons, recognition and human judgement. Haines et al. include human evaluation, while image-generation work also uses text and writer recognition to assess the output~\cite[pp.~12--14]{haines2016},~\cite[pp.~8--9]{kang2020}. For the present robot, model scores help choose a candidate before drawing, but final physical output also needs its own read-back and geometric assessment. A classifier used to select the candidate is not an independent human judgement of its style.

\subsection{Physical reproduction and pen-path mapping}
There are three relevant families of mechanism for reproducing text on paper: Cartesian systems, serial-chain robots and parallel-chain robots. A Cartesian system consists of a gantry with independent $x$--$y$ movement. Serial-chain robots consist of linkages connected in a chain, like a typical robot arm. Parallel mechanisms use several linked chains to support the moving platform~\cite[pp.~5--8]{teikari2016}. These arrangements differ in their workspace, mechanical structure and coordinate transformations.

A Cartesian gantry is suitable here because the page is flat and the main motion requirement is two-dimensional. Each axis can be commanded and calibrated separately, with a further pen-lift mechanism for travel. This makes the conversion from an image-derived path to physical coordinates direct, while leaving frame alignment, transmission compliance and missed steps as practical accuracy concerns.

A pen path is stored as a list of points, and a recovered stroke usually has more points than a plotter needs. Polyline simplification reduces this list while retaining the stroke's shape. The Douglas--Peucker method joins the endpoints, finds the point furthest from that line and retains it when its distance exceeds a chosen tolerance, then repeats on the resulting parts~\cite{douglas1973}. A small tolerance follows the stroke closely; a large tolerance removes more points but can cut corners. When handwriting exists only as an image, the path must first be recovered, for example by thinning the ink to a centreline~\cite[p.~237]{zhang1984}.

These operations are needed to turn the project's stored glyphs into manageable G-code trajectories. Simplification controls geometric detail, while the thinning and tracing determine which strokes are available to the synthesiser. Neither recovers the original writing speed or pressure from a photograph, so the controller must define its own drawing feed and pen transitions.

There are multiple methods of linear actuation that could be considered for a Cartesian plotting system, including belt drives, rack-and-pinion drives and screw drives~\cite{iqs2024}. The project uses timing belts and pulleys to convert motor rotation into movement along supported rails. This choice prioritises a straightforward, affordable page-sized mechanism with sufficient travel speed; its nominal step resolution must still be distinguished from measured positioning accuracy.

\subsection{Application of the literature}
\label{sec:lit-application}
\label{sec:lit-embedded}
The literature supports a system built around offline pages, complete-line recognition and explicit glyph geometry. Image conditioning and curved line separation prepare photographed input; convolutional features support separate text and writer outputs; transcript-guided extraction turns observed handwriting into reusable paths. A glyph library then supplies controllable geometry without training an additional handwriting generator. The Cartesian gantry provides the physical output stage. Section~\ref{sec:approach} explains how these choices were combined and developed, while Section~\ref{sec:design} gives their calculations, implementation and verification.
